% !TeX root = ../main.tex
\section{Suggestions for Future Work}

The future directions of this project follow directly from the main limitations identified in the discussion. Although the thesis demonstrated that embedded soccer analytics is feasible on the Jetson Orin Nano, the evaluation was limited to a single hardware platform, a selected set of broadcast-style datasets, a constrained set of full-stack statistics, and a partial exploration of the broader compression design space. In addition, the final end-to-end system achieved a practically useful operating point, but did not fully reach the most ambitious real-time throughput target. Future work should therefore not be viewed simply as feature expansion, but as a set of focused research investigations aimed at improving generalisability, deployment reliability, and end-to-end analytical value.

\subsection{Extending Full-Stack Match Understanding}

One important direction is to extend the system beyond the current set of outputs toward richer match statistics and event understanding. In particular, future work could investigate pass completion percentage, turnovers, ball progression sequences, possession chains, and more detailed player movement summaries. This would make the system more useful to coaches, analysts, and broadcasters, and would help bridge the gap between low-level detections and higher-level football understanding.

However, this extension is not a straightforward implementation task. Reliable pass completion or turnover analysis depends on stable ball detection, robust player tracking, accurate team assignment, and temporally consistent possession reasoning. Errors in any one stage can propagate to the final statistic, especially during occlusion, rapid camera motion, or ambiguous multi-player interactions. Future studies should therefore examine temporal modelling strategies, confidence-aware event rules, and methods for validating event-level outputs against manually annotated sequences. The impact of this work would be significant, since it would determine whether embedded systems of this type can move beyond demonstration-level analytics and provide outputs that are useful in realistic tactical or broadcast settings.

\subsection{System-Level Optimisation for Real-Time Deployment}

A second major area for further investigation is end-to-end system optimisation. The results showed that model compression alone was not sufficient to close the remaining gap to the target real-time operating point, which suggests that pipeline-level design is now a central research problem. Future work should evaluate asynchronous execution between inference and post-processing, adaptive object-detection intervals, less frequent team-update scheduling, and hybrid approaches in which some pitch or geometric operations are replaced with lighter classical methods \cite{okuma2004automatic}. Another promising direction would be to study dynamic model selection, where smaller or faster models are used under high-load conditions and more accurate models are used when the scene is less demanding.

The main difficulty here is that throughput improvements may come at the cost of stale detections, delayed state updates, or unstable inter-stage synchronisation. A faster system is not necessarily a better system if the resulting statistics become temporally inconsistent or analytically unreliable. Future research should therefore measure not only FPS, but also the statistical effect of scheduling decisions on possession estimation, pass detection, and tracking continuity. This line of work is important because it is likely to determine whether the prototype can become a genuinely real-time embedded application rather than a near-real-time research demonstrator.

\subsection{Advanced Compression and Training Strategies}

The compression study in this thesis established broad trends, but several important directions remain open. Future work should investigate quantization-aware training in addition to post-training quantization, since the results suggest that INT8 deployment offers strong efficiency gains but can introduce substantial accuracy loss. It would also be valuable to compare structured and unstructured pruning more systematically, explore sparsity-aware retraining, and test a wider range of teacher-student pairings for knowledge distillation. The preliminary input-resolution experiments could also be extended by training and evaluating models under controlled resolution mismatch, to determine whether low-cost training or inference configurations can retain acceptable robustness.

These studies are likely to be challenging because compression behaviour is highly task-dependent. The results already showed that pitch keypoint models are much more sensitive than object detectors to aggressive pruning and reduced numerical precision, so future work will need careful experimental design to avoid drawing conclusions from unstable or unfair comparisons. There is also a practical cost in terms of training time, hardware usage, and calibration effort when many compression variants are evaluated. Despite these difficulties, the impact of such work would be substantial: a better understanding of compression-aware training could improve the accuracy--efficiency frontier of embedded sports analytics and produce deployment candidates that are both faster and more trustworthy.

\subsection{Robustness, Generalisation, and Cross-Platform Validation}

Another key limitation of the present work is that the evaluation was carried out under relatively controlled conditions. Future research should therefore examine how well the trained and compressed models generalise to lower-quality footage, non-broadcast camera placements, different lighting conditions, amateur or youth-level matches, and more diverse pitch environments. In parallel, the full-stack system should be tested over longer-duration runs to evaluate thermal stability, sustained throughput, and failure modes under extended operation. Cross-platform benchmarking on additional edge devices would also be valuable, since deployment trade-offs may differ considerably across CPU, GPU, and memory configurations.

The main challenge in this direction is that robust evaluation requires new annotated data, careful licensing and collection procedures, and access to multiple hardware platforms. Domain shift may also make it difficult to determine whether performance degradation is caused primarily by the model, the compression method, or the broader system assumptions. Even so, this is one of the most important areas for future work, because it determines whether the conclusions of this thesis remain valid outside the specific datasets and hardware used here. Demonstrating broader robustness would strengthen both the scientific contribution and the practical impact of the work, especially for lower-budget organisations that may operate in less controlled recording conditions than professional broadcast environments.

\subsection{Research Value of Continued Development}

Taken together, these future directions show that the next stage of this project should focus on transforming a strong prototype into a more general and dependable embedded analytics framework. The importance of carrying out this work lies not only in improving one specific soccer-analysis system, but in advancing the broader understanding of how modern vision pipelines can be compressed, validated, and deployed responsibly on constrained hardware. Further research in these areas could benefit both the sports analytics domain and the wider field of embedded AI, where the same questions of efficiency, robustness, and practical usefulness arise across many real-world applications.
